% 1 Statement of the problem (all; first draft: Aryan). Sources: deck slides 2-6, claims.csv rows S1, S1/O1, S5/O3.
\section{Statement of the problem}\label{sec:problem}

Artificial intelligence is usually discussed as software, but it has a physical body. Every model is trained and run
in a data centre: a building packed with servers whose chips draw electricity, shed heat that has to be removed, often
by evaporating water, and are made of copper, gallium, cobalt and rare earths that are mined in a few places and end up
as electronic waste. In this report ``AI'' means that infrastructure, not the algorithms.

Data centres used about 415\,TWh in 2024, roughly 1.5\% of the world's electricity, and the International Energy
Agency's base case has them reaching about 945\,TWh by 2030 \citep{iea2025}. That growth lands on an Earth system that
is already outside most of its safe operating space. Each assessment of the planetary boundaries has found more of them
crossed: three of seven in 2009 \citep{rockstrom2009}, four in 2015 \citep{steffen2015pb}, six of nine in 2023
\citep{richardson2023} and seven of nine in the latest Planetary Health Check \citep{sakschewski2025,src_pb}. Data
centres press on several of them at once: climate through their electricity, freshwater through cooling and power-plant
water, and novel entities and biosphere integrity through e-waste and mining.

A share of 1.5\% sounds small, and globally it is. But it is not spread evenly. The United States (45\%), China (25\%)
and Europe (15\%) account for 85\% of data-centre electricity \citep{iea2025}, and the large AI sites tracked by Epoch
AI are more concentrated still: 83\% of their 15\,GW is in the United States \citep{epoch2026dc}. The impacts
therefore fall on particular grids, river basins and communities. A recent global study placed one quarter of the
world's data centres in water-stressed areas \citep{rosa2026}, but it counts facilities: a 10\,MW server room and a
2\,GW AI campus count the same, and their water use is not estimated. No assessment we know of weighs data centres by
size and relates them to the water and the people around them. That is the gap this project addresses.

\paragraph{Three phases and one puzzle} We frame the history in three phases since 1950. In the \emph{Great
Acceleration} (1950--2010) energy use, \cotwo\ and material extraction rose together \citep{steffen2015ga}. In the
\emph{cloud era} (2010--2018) computing grew enormously, yet global data-centre electricity barely moved: by one widely
used estimate, compute instances rose by about 550\% while energy use rose by about 6\% \citep{masanet2020}. In the
\emph{AI era} (2019 onward) that decoupling broke: US data-centre electricity rose from 76 to 176\,TWh between 2018
and 2023 \citep{shehabi2024}. Our organising question is therefore: \emph{why did cloud growth leave energy use almost
flat, while AI growth does not?}

\begin{finding}{Terms used in this report}
\textbf{MW and TWh.} Power (MW, GW) is a rate; energy (TWh) is power over time. 1\,GW running all year is about
8.8\,TWh. Epoch AI reports \emph{facility} power, the whole building, not just the computers.
\textbf{PUE} (power usage effectiveness) = facility energy $\div$ IT energy; 1.5 means an extra 50\% goes to cooling
and losses. \textbf{WUE} (water usage effectiveness) = litres of water consumed on site per kWh of IT energy.
\textbf{Grid water intensity} = litres consumed at power plants per kWh generated. \textbf{Withdrawal vs
consumption:} water withdrawn is taken from a river or aquifer, and part is returned; water consumed is evaporated and
lost to the basin. \textbf{Water stress} (WRI Aqueduct) = annual withdrawals as a share of a basin's renewable supply:
above 40\% is \emph{high}, above 80\% \emph{extremely high}.
\end{finding}
